\documentclass[11pt]{article}

\newcommand{\cnum}{Math 245B}
\newcommand{\ced}{Winter 2026}
\newcommand{\ctitle}[4]{\title{\vspace{-0.5in}\cnum, \ced\\Assignment #1: #2}\author{\vspace{-0.35in}\\#3, #4}}
\usepackage{enumitem}
\usepackage[usenames,dvipsnames,svgnames,table,hyperref]{xcolor}
\usepackage{amsmath, amsfonts}
\usepackage{graphicx} % Required for inserting images
\usepackage{float}
\usepackage{listings}
\usepackage{xcolor}
\usepackage{hyperref}
\usepackage{comment}

\renewcommand*{\theenumi}{\alph{enumi}}
\renewcommand*\labelenumi{(\theenumi)}
\renewcommand*{\theenumii}{\roman{enumii}}
\renewcommand*\labelenumii{\theenumii.}
\newcommand{\notiff}{%
  \mathrel{{\ooalign{\hidewidth$\not\phantom{"}$\hidewidth\cr$\iff$}}}}

\definecolor{codegreen}{rgb}{0,0.6,0}
\definecolor{codegray}{rgb}{0.5,0.5,0.5}
\definecolor{codepurple}{rgb}{0.58,0,0.82}
\definecolor{backcolour}{rgb}{0.95,0.95,0.92}
\definecolor{header}{HTML}{205459}
\definecolor{solution}{HTML}{204d52}

\newcommand{\solution}[1]{{{\textcolor{header}{Solution:} \textcolor{solution}{#1}}}}
\setlist[enumerate]{label=(\alph*)}

\lstdefinestyle{mystyle}{
    backgroundcolor=\color{backcolour},
    commentstyle=\color{codegreen},
    keywordstyle=\color{magenta},
    numberstyle=\tiny\color{codegray},
    stringstyle=\color{codepurple},
    basicstyle=\ttfamily\footnotesize,
    breakatwhitespace=false,
    breaklines=true,
    captionpos=b,
    keepspaces=true,
    numbers=left,
    numbersep=5pt,
    showspaces=false,
    showstringspaces=false,
    showtabs=false,
    tabsize=2
}


\begin{document}
\ctitle{\#1}{}{Asher Christian}{006-150-286}
\date{}
\maketitle

\section{Monday}
Strong Markov Property
\[
    1_{\tau  <\infty}   \mathbb{E}[Y \circ \theta_\tau | \mathcal{F}_\tau] = 1_{\tau < \infty}\mathbb{E}_{X_\tau} Y
\] 
\[
    \mathcal{F}_\tau = \{A : A \cap \{\tau = n\} \in \mathcal{F}_n\}
\] 
since it is a stopping time we have $\{\tau = n\} \in \mathcal{F}_n$
\[
    = \{A : A \cap \{\tau = n\} \in \mathcal{F}_n \cap \{\tau = n\}\}
\] 
\[
    = \{A  \in \bigcup_{n=1}^{\infty} \left(\mathcal{F}_n \cap \{\tau = n\}\right)\}
\] 
e.g.
\[
A = \bigcup_{n} \left\{X_{n-1} = 5\} \cap \{\tau = n\}\right)
\] 
r.h.s = $g(X_\tau)$ 
$g(a) = \mathbb{E}(Y | X_0 = a)$
Example:
simple random walk half half from $X_0 = 0$
\[
    \tau = \min\{n  > 0 : X_n \in \{-5,10\}\}
\] 
what is
\[
    \mathbb{P}(X_\tau + 2 = 12) = \mathbb{E}[1_{X_{\tau + 2} = 12}] = \mathbb{E}(1_{X_2 = 12} \circ \theta_\tau)
\] 
\[
    = \mathbb{E}[ 1_{\tau < \infty} \left(1_{X_2 = 12} \circ \theta_\tau\right)]
\] 
\[
    = \mathbb{E}\left( \mathbb{E}[1_{\tau < \infty} ( 1_{X_2 = 12} \circ \theta_\tau) | \mathcal{F}_{\tau}] \right)
\] 
\[
    = \mathbb{E}\left( 1_{\tau < \infty} \mathbb{E}[1_{X_2 = 12} \circ \theta_\tau| \mathcal{F}_\tau]\right) 
\] 
\[
    = \mathbb{E}\left(1_{\tau < \infty} \mathbb{E}_{X_\tau} 1_{X_2 = 12}\right)
\] 
\[
    = \mathbb{E}[g(X_{\tau})]
\] 
where
\[
    g(a) = \mathbb{P}(X_{12} = 12 | X_0 = a)
\] 
\[
    = \mathbb{P}(X_{\tau} = -5) g(-5) + \mathbb{P}(X_{\tau} = 10) g(10)
\] 
\[
= \frac{1}{3} g(10) = \frac{1}{12}
\] 
Second example
THMC $(S, \mathcal{S})$
Def: 
 \[
     \rho_{xy} : \mathbb{P}(\exists n > 0 , X_n = y | X_0 = x)
\] 
\[
    = \mathbb{E}[1_{\exists n > 0, X_n  =y} | X_0 = x]
\] 
\[
    \mathbb{P}(\exists n_1 > n_2 > 0 , X_{n_1} = z, X_{n_2} = y | X_0 = x) = \mathbb{P}(A | X_0 = x)
\] 
what is 
\[
    \mathbb{E}[1_A | X_0 = x] = \rho_{xy} \rho_{yz}
\] 
\[
    1_A \sim Y \circ \theta_\tau
\] 
\[
    \tau = T_y \quad Y = 1_{\exists n > 0, X_n = z}
\] 
then
\[
    Y \circ \theta_\tau = 1_{\exists n > \tau , X_n = z}
\] 
\[
    1_A \sim Y \circ \theta_\tau 1_{\tau < \infty}
\] 

\[
    \mathbb{E}_x (1_{\tau < \infty} Y \circ \theta_\tau)
\] 
\[
    = \mathbb{E}_x(1_{\tau < \infty} \mathbb{E}(Y \circ \theta_\tau | \mathcal{F}_\tau]) = \mathbb{E}_x(1_{\tau < \infty} g(X_\tau))
\] 
where $g(a) = \mathbb{E}[Y | X_0 = a] = \mathbb{P}(\exists n > 0, X_n = z | X_0 = a)$ 
\[
    = \mathbb{E}_x(1_{\tau < \infty} \rho_{yz}) = \rho_{yz} \rho_{xy}
\] 
strong strong markov property
\[
    1_{\tau < \infty} \mathbb{E}[Y_{\tau} \circ \theta_{\tau} | \mathcal{F}_{\tau}] = \mathbb{E}_{X_\tau}(Y_\tau) 1_{\tau < \infty}
\] 
for example $X_{2\tau} = [X_{\tau} \circ \theta_\tau]$
for the rhs
\[
    = g(X_{\tau} , \tau)
\] 
\[
    g(a,m) = \mathbb{E}[Y_m | X_0 = a]
\] 
Random walk
\[
X_n = \sum_{k=1}^{n}\xi_k 
\] 
i.i.d
\[
\mathbb{E} \xi_k = 0
\] 
\[
\mathbb{P}(\xi_k > a) = \mathbb{P}(\xi_k < -a)
\] 
for any $a > 0$ and $ \forall n $
 \[
\mathbb{P}(\exists m , 0 \le m \le n, X_m \ge a) \le 2 \mathbb{P}(X_n \ge a)
\] 
\section{Wednesday}
Application of SMCP
Example $X_n = \sum_{k=1}^{n}\xi_k$ with $\xi_k$ i.i.d  $\xi_k \sim \mu$ symmetric
\[
\forall a, \mathbb{P}(\exists m, 0 \le m \le N, X_m \ge a) \le 2 \mathbb{P}(X_N \ge a)
\] 
\[
    \tau - \min\{n > 0, X_n \ge a\}
\] 
Proof:
\[
    \mathbb{P}(X_N \ge a) \ge \mathbb{P}(\tau \le N,  X_N - X_{\tau} \ge 0) = \mathbb{P}(B)
\] 
%are they not equal
$B = \bigcup B_m$ where $B_m = \{\tau = m\} \cap \{X_n \ge X_m\}$ clearly $B_m \subset A$ so $B \subset A$.
\[
    = \mathbb{E}[1_{\tau \le N} 1_{X_N - X_{tau} \ge 0}]
\] 
For strong markov chain property we need to write it into $Y_\tau \circ \theta_\tau$\\
We need to define $Y_m$ first
 \[
     Y_m = 1_{X_{N-m} - X_0 \ge 0}
\] 
for $m \le N$
Then
\[
    Y_\tau \circ \theta_\tau = 1_{X_N - X_{\tau} \ge 0}
\] 
and we know
\[
\mathbb{P}(Y_m = 1) \ge \frac{1}{2} 
\] 
\[
    = \mathbb{E}[1_{\tau \le N} \mathbb{E}[ Y_{\tau} \circ \theta_\tau | \mathcal{F}_\tau]]
\] 
\[
    = \mathbb{E}[1_{\tau \le N} g(X_{\tau},\tau)]
\] 
\[
    g(X_\tau,\tau) = \mathbb{E}_{X_\tau}Y_\tau
\] 
\[
    g(a,m) = \mathbb{E}[Y_m | X_0 = a] = \mathbb{E}[Y_m] \ge \frac{1}{2}
\] 
so
\[
    \ge \frac{1}{2} \mathbb{E}[1_{\tau \le N}] = \frac{1}{2} \mathbb{P}(\tau \le N)
\] 
simple random walk $N = 10^{4}$ we can calculate $X_N$ distribution and
$ \mathbb{P}(X_N \ge 2\cdot10^{2})$ 
what about
\[
    \mathbb{P}(\max_{m \le n} X_m \ge 2 \cdot 10^{3}) \le 2 \mathbb{P}(X_N \ge 2 \cdot 10^{3})
\] 
Recurrent and transient
if $\rho_{xx} = 1$ then $\rho_{yx} = 1$ if $x \to y$ since if  $\rho_{yx} < 1$ then there
is a finite probability you got to $y$ and dont come back so there is a finite probability you never come back so  $\rho_{xx} < 1$\\
Let $S$ be countable
Definition: recurrent $\rho_{xx} = 1$ and transient $\rho_{xx} < 1$
\[
    \rho_{xx}^{(2)} = \rho_{xx} \rho_{xx} = \rho_{xx}^2
\] 
\[
    \rho_{x\to y\to z} = \rho_{xy}\rho_{yz}
\] 
similarly
\[
    \rho_{xx}^{(m)} = \rho_{xx}^{m}
\] 
Def
\[
    N_x = |\{n > 0 : X_n = x\}|
\] 
\[
    \mathbb{P}_x(N_x \ge 1) = \rho_{xx}
\] 
similarly
\[
    \mathbb{P}_x(N_x \ge 2) = \rho_{xx}^2
\] 
\[
    \mathbb{E} N_x = \sum_{k}^{} \mathbb{P}(N_x \ge k) = \frac{\rho_{xx}}{1- \rho_{xx}} = 
    \begin{cases}
        \infty & R\\
        < \infty & T
    \end{cases}
\] 
\[
    \mathbb{E}_x N_x = \sum_{k}^{} \mathbb{E}_x 1_{X_k = x} = \sum_{k}^{} \mathbb{P}_x(X_k = x)
\] 

\section{Friday}

Lemma: 
\[
    \mathbb{E}_x N_x = \frac{\rho_{xx}}{1 - \rho_{xx}} = \sum_{k=0}^{\infty} \mathbb{P}_x(X_k = x)
\] 
This can show whether a state is recurrent or transient.
Th example is simple random walk in $ \mathbb{Z}^{d}$ we can show that
\[
\mathbb{P}_x(X_2k = x) \sim k^{-\frac{d}{2}}
\] 
so
\[
    \rho_{xx} < 1 \quad d \ge 3 \qquad \rho_{xx} = 1 \quad d \le 2
\] 
Lemma:
if $x \to y$  and  $y \to x$ where we say  $x \to y \iff \rho_{xy} > 0$
and $\rho_{xx} = 1 \implies \rho_{yy} = 1$
$x \sim y \iff x \to y, y \to x$ is an equivalence relation.
Proof:
suppose  $x \to y$ so  $\exists n$ such that  $p^{(n)}_{xy} > 0$ 
and since $y \to x$ there exists  $m$ such that  $p_{yx}^{(m)} > 0$ we want to prove
\[
\mathbb{E}_y N_y = +\infty = \sum_{k=1}^{\infty} \mathbb{P}_y(X_k = y) = +\infty
\] 
\[
    = \sum_{k=1}^{\infty} p_{yy}^{(k)} = +\infty
\] 
where
\[
    p_{xy}^{(n)} = \mathbb{P}(X_n = y| X_0 = x)
\] 
\[
    \sum_{k=1}^{\infty} p_{yy}^{(k)} \ge \sum_{k=1}^{\infty}p^{(k+m+n)}_{yy} \ge \sum_{k=1}^{\infty}p_{yx}^{(m)}p^{(k)}_{xx}p^{(n)}_{xy} \ge C \sum_{k=1}^{\infty}p_{xx}^{(k)} 
\] 
so if the rhs is infinite then the left hand side is infinite.
\section{Lemma}
If $\rho_{xx} = 1$ and $x \to y$ then  $\rho_{yy} = 1$ and $\rho_{yx} = 1$.
Let
\[
    T = \min\{n > 0, X_n = y , 1 \le m < n, X_m \ne x\}
\] 
\[
\mathbb{P}(T_x < \infty | X_0 = x) = \mathbb{P}(T_x < \infty , T> \infty | X_0 = x) + \mathbb{P}(T_x < \infty, T < \infty | X_0 = x)
\] 
%shouldnt we say show that P(T < inft) > 0
\[
 = \mathbb{P}_x(T > \infty) + \mathbb{P}(T_x < \infty | T< \infty , X_0 = x)  \mathbb{P}(T < \infty)
\] 
\[
   \implies 1 =  \mathbb{E}[1_{T_x < \infty} | T< \infty X_0 =x]
\] 
\[
    = \mathbb{E}[1_{T_x < \infty} 1_{T < \infty} | T< \infty X_0 = x]
\] 
\[
    = \mathbb{E}[1_{T < \infty} \circ \theta_T
\] 
\[
    \sim  \mathbb{E}_{X_T} T_x 1_{T< \infty}
\] 
and since $X_T = y$ we have equal  $\rho_{yx}$\\
\\
If $|S| < \infty$ then there exists $x \in X$ such that  $\rho_{xx} = 1$
Proof: If $x \to y$  then
\[
\begin{cases}
    \rho_{yy} = 1 & \mathbb{E}_x N_y = +\infty\\
    \rho_{yy} < 1 & \mathbb{E}_x N_y < \infty
\end{cases}
\] 
And since $|S| < \infty$
\[
\sum_{x \in S}^{} N_x = +\infty
\] 
so
\[
    \mathbb{E} \sum_{x \in S}^{} N_x = +\infty \implies \sum_{x \in S}^{} \mathbb{E}_{*}N_x = +\infty 
\] 
so for all $x \in S$ there exists y such that  $ \mathbb{E}_x N_y = \infty$\\
If $x \to y$ and  $y \not\to x$ then  $x $ is transient




\end{document}
